% =====================================================================
%  Partie 0 — Introduction générale
% =====================================================================
\part{Introduction générale}

\chapter{Introduction et objectifs}

\section{Présentation du cours}
Ce document est le support de cours de l'enseignement A 567 Analyse de cycle de vie, dispensé à l'ENSAP Lille. Ce document est pensé pour être un récapitulatif des informations essentielles présentées et discutées en cours : il ne se substitue pas au cours ni aux TD d'application, il en est plutôt un résumé synthétique.

\textbf{Ce document est en cours d'élaboration et fait suite à la lutte de 2026 dans les ENSA, il n'est pas encore finalisé. Une première version devrait arriver en décembre 2026 : il peut donc contenir des erreurs. Si vous repérez des erreurs, merci de me les faire remonter.}

Dernière mise à jour : \today.\\

Ce cours est composé de 4 parties :
\begin{itemize}
  \item \textbf{Partie 1 : Introduction générale.} Présentation du cours et des objectifs.
  \item \textbf{Partie 2 : Contexte environnemental.} Présentation des enjeux environnementaux, de l'impact du secteur du bâtiment et du cadre réglementaire actuel.
  \item \textbf{Partie 3 : Méthode de l'ACV.} Présentation de l'historique et des principes de l'analyse de cycle de vie, ainsi que son fonctionnement (indicateurs et modules).
  \item \textbf{Partie 4 : ACV et architecture.} Présentation de l'utilisation de la méthode d'ACV en architecture, à travers les données et outils disponibles.
\end{itemize}

\section{Objectifs pédagogiques}

Les objectifs de ce cours sont de :
\begin{itemize}
  \item Comprendre les principes de la méthode d'Analyse de Cycle de Vie (ACV) appliquée à l'architecture.
  \item Utiliser l'ACV comme un outil d'aide à la conception dès les premières phases du projet.
  \item Identifier et comparer les impacts environnementaux de différentes solutions architecturales et constructives.
  \item Développer un regard critique sur les résultats d'une ACV pour éclairer les choix de projet.
\end{itemize}

\section{Organisation}

Le cours alterne cours magistraux et travaux dirigés d'application.

\section{Références}

La bibliographie complète est présente en fin de document. 
Ce cours est principalement inspiré des cours d'analyse de cycle de vie de l'École des Ponts et Chaussées, par Adélaïde Féraille, et de l'ouvrage \textit{Analyse du cycle de vie : Comprendre et réaliser un écobilan}\footfullcite{saade_analyse_nodate} de Myriam Saadé et Olivier Jolliet. \\

Voici une liste de ressources utiles à consulter dans différents formats :


\begin{encadre}{Sites web}
  \begin{itemize}
  \item Base INIES, \url{https://www.base-inies.fr}
  \item Unisson(s), boîte à outils, \url{https://unissons.eu/boite-a-outils/}
  \end{itemize}
\end{encadre}

\begin{encadre}{Bandes dessinées}
  \begin{itemize}
    \item \textit{Le Monde sans fin}, Jancovici et Blain (2021)
    \item \textit{Ressources : un défi pour l'humanité}, Bihouix et Perriot (2024)
  \end{itemize}
\end{encadre}

\begin{encadre}{Vidéos}
  \begin{itemize}
    \item \textit{Écolo mon cul !}, Arte, 8 vidéos de 10 min, \url{https://www.youtube.com/playlist?list=PLFqwNouKXOkE}
  \end{itemize}
\end{encadre}
